%==============================================================================
\section{Computational study}\label{sec:comp}
%==============================================================================

Computation here is how the theory is tested. The quantity the
structural results single out, the coverage bound $q$, is also the root relaxation of
every configuration-based method, so a benchmark that records optima and bounds
measures the theory's central variable directly. The campaign also settles a question
that predates the theory and motivated the Benders solver in the first place: the SSP
has a polynomially solvable loading subproblem, which is the textbook precondition for
Benders decomposition, and whether a decomposition can exploit it against modern
compact models had not been tested.

\paragraph{What the section asks, and in what order.}
Each subsection states the question it answers before the measurement that answers it.
Section~\ref{ssec:bench} fixes the protocol and Section~\ref{ssec:validation} asks
whether the implementations are sound. Section~\ref{ssec:results-solved} asks how far
each method gets; Section~\ref{ssec:results-bbc} whether the difficulty is finding good
solutions or proving them optimal, and whether the standard strengthenings move it;
Section~\ref{ssec:results-tight} what separates the instances a method closes from
those it does not. Section~\ref{ssec:results-bnp} puts the same question to the
branch-and-price prototypes, which report their relaxation value directly and so answer
it without reference to solving times. Section~\ref{sec:disc} then asks why the answer
is what it is.

\subsection{Experimental design}\label{ssec:bench}

\paragraph{Instances.}
The benchmark uses the standard SSP instance families, taken from the collection
assembled by \citet{Otiai2024}: the sets of \citet{Catanzaro2015}, of
\citet{Crama1994}, and four of the sets of \citet{Laporte2004}. Every instance in
every family is run; none is skipped, and no instance is selected on the basis of an
outcome. Table~\ref{tab:families} gives the sizes.

\begin{table}[h]\centering
\footnotesize
\begin{tabular}{@{}lrrrl@{}}
\toprule
Family & Instances & Jobs $n$ & Tools $\lvert T\rvert$ & Source \\
\midrule
Catanzaro & $171$ & $8$--$40$  & $5$--$60$  & \citet{Catanzaro2015} \\
Crama     & $160$ & $10$--$40$ & $10$--$60$ & \citet{Crama1994} \\
Laporte3  & $340$ & $8$        & $15$--$25$ & \citet{Laporte2004} \\
Laporte4  & $330$ & $9$        & $15$--$25$ & \citet{Laporte2004} \\
Laporte5  & $340$ & $15$       & $15$--$25$ & \citet{Laporte2004} \\
Laporte7  & $80$  & $10$--$15$ & $10$--$20$ & \citet{Laporte2004} \\
\midrule
Total     & $1{,}421$ & & & \\
\bottomrule
\end{tabular}
\caption{The instance families. An instance is a matrix together with a magazine
capacity. The Crama collection publishes the same matrices at four different
capacities, so its $40$ matrices give $160$ instances; the same convention is applied
throughout.}
\label{tab:families}
\end{table}

\paragraph{Protocol.}
Every run of the Benders solver and of the compact baselines is given a single uniform
time limit of one hour and $16$ gigabytes of memory. There is no early stopping and no
per-configuration retirement rule. All other solver parameters are left at their
defaults, so that the comparison measures the formulations rather than parameter
tuning. The branch-and-price prototypes are run separately, under a ten-minute limit,
for the reason given in Section~\ref{ssec:results-bnp}.

\paragraph{Hardware and software.}
All runs were executed on the LIMOS SLURM cluster, confined to a homogeneous pool of
identical nodes (AMD EPYC~7452, $2\times32$ cores, $512$\,GB of memory), so that
wall-clock times are comparable across runs. The mixed-integer solver is CPLEX~22.1.1
throughout, pinned to four dedicated threads per run. The branch-and-price prototypes
run single-threaded under PySCIPOpt. The campaign was distributed over $120$ SLURM
array shards and produced $16{,}994$ individual runs.

\paragraph{Configurations.}
Twelve solver configurations were run. Three are the reimplemented compact baselines:
F4 of \citet{Catanzaro2015}, written \texttt{CATZ-F4}; LSS of \citet{Laporte2004}; and
SSPMF of \citet{daSilva2024}. F4 was chosen over the F5 that \citet{Catanzaro2015}
recommend because it is the smaller of the two models that share the F3 relaxation, and
implementing F5 did not fit in the time available. The compact side of this comparison
is therefore weaker than the literature allows, and every statement made from it is
made in that knowledge; Section~\ref{ssec:results-summary} reports what follows, which
is that the baselines win anyway. The remaining nine are the Benders solver, in an ablation
over its cut families and its strengthenings. The naming is as follows:
\texttt{BBC-LP} uses the dual optimality cuts, \texttt{BBC-K} the combinatorial cuts,
\texttt{+T} adds the triplet bounds, \texttt{+F} enables fractional separation,
\texttt{+C} the conflict-graph row, \texttt{+H} the hybrid-genetic incumbent,
\texttt{+P} the Pareto-optimal cut lifting, and \texttt{+ACC} all three strengthenings
together.

\paragraph{The reported quantity.}
The common metric is the empty-start cost of the sequence a method returns, evaluated
by the loading oracle. Each formulation has its own native objective, and two of them
use the free-initial convention, so the native values are not comparable across
methods. The empty-start cost of the returned sequence is defined for every method
regardless of its formulation, and it is what makes the cross-method comparison of
Section~\ref{ssec:validation} meaningful.

\subsection{Validation}\label{ssec:validation}

The implementations were subjected to independent computational checks before any
result was recorded, and again from the results themselves. No amount of such checking
proves correctness; what follows is what was checked and what it found.

\paragraph{Before the campaign.}
Every exact solver was checked against brute-force optima on a bank of small instances,
in every combination of its option flags. The loading oracle itself was verified
against an exact dynamic program over magazine states; the verification described in
Appendix~\ref{app:verify} repeats this on $349{,}872$ instance-and-sequence pairs and
finds no disagreement.

\paragraph{Identifying an instance.}
An instance is identified by the family it comes from, its file name, its number of
jobs, its number of tools, and its magazine capacity. All five are needed. The Crama
collection reuses the same file names at four different capacities, so two runs that
agree on the name can be runs on different instances. Keying on the name alone merges
rows describing different instances and produces apparent disagreements between methods
where none exists. The comparison below uses the full identifier.

\paragraph{Cross-method agreement.}
Of the $1{,}421$ instances, $1{,}115$ were solved to proven optimality by at least one
method and $999$ by at least two. On those $999$ there are $41{,}673$ pairs of methods
that both proved an optimum. \emph{In none of them do the two disagree.} Every pair of
methods that both closed an instance returned a sequence of the same empty-start cost.

\paragraph{The two conventions, measured.}
Proposition~\ref{prop:conv} states that the empty-start and free-initial costs of a
sequence differ by the size of the initial load. The campaign confirms it on real
instances. The SSPMF model reports the free-initial cost natively, and on all
$1{,}038$ instances it closed, the difference between its reported value and the
empty-start cost of the same sequence equals the magazine capacity $b$ on $1{,}029$ of
them. On the other nine it is smaller than $b$, and those are exactly the instances on
which an optimal sequence begins with a job that needs fewer than $b$ tools, so the
magazine is not filled at the start. The two baselines that report the empty-start cost
natively, \texttt{CATZ-F4} and LSS, show a difference of zero on all $1{,}729$ runs in
which they closed an instance.

\paragraph{Tools that no job needs.}
Four matrices in the collection contain a tool that appears in no job's requirement
set: Crama \texttt{s2n007} and \texttt{s2n009}, Laporte3 \texttt{L28-7} and Laporte5
\texttt{L6-8}. Counted with the capacities at which they appear, that is ten of the
$1{,}421$ instances. The coverage bound of Corollary~\ref{cor:lb} must therefore be read as
$q=\lvert\U\rvert-b$, using the set $\U$ of tools that some job actually requires, and
not as $\lvert T\rvert-b$ using the tool count declared in the instance file. On those
ten the two differ, and the smaller value is the correct one: on Laporte3
instance \texttt{L28-7}, which declares $25$ tools of which one is never required and
has $b=5$, the optimum in the free-initial convention is $19$, which is
$\lvert\U\rvert-b=24-5$ and not $\lvert T\rvert-b=20$. All coverage-bound figures in
this section use $\lvert\U\rvert$.

\paragraph{A defect found and corrected.}
The Benders master carries the coverage row $\theta\ge\lvert\U\rvert$ of
Section~\ref{ssec:bbc} as an explicit constraint. It has to: unlike every other method
here, the master has no formulation of the loading problem from which that bound would
follow, so without the row it is not available to the relaxation at all.
Observation~\ref{obs:rootq} shows how much work the row does --- on $87.4\%$ of runs the
root relaxation is exactly $\lvert\U\rvert$ and nothing else contributes.

\paragraph{Runs that failed.}
Of the $16{,}994$ runs, $427$ returned no result. Forty-nine are a failure of the
harness itself while closing its log file on the cluster filesystem, and the remaining
$378$ are crashes of the worker subprocess. They are not spread evenly: $352$ of the
crashes belong to the LSS baseline, $239$ of those to the Laporte5 family.

The two groups have different causes, and one of them is recorded.

\emph{The Benders failures.} Of the $75$ failures scattered across the nine Benders
configurations, at most $26$ in any one, $49$ carry an explicit exception:
\texttt{OSError: [Errno 127] Key has expired}. That is expiry of the job's filesystem
credentials on a cluster whose home directories are served over an authenticated mount.
The run lost the ability to write its own output; nothing about the instance or the
solver is implicated. These are infrastructure losses, and they change no comparison
drawn between configurations.

\emph{The LSS crashes.} The remaining $352$ are recorded as
\texttt{subprocess crash} with an empty message: the worker was terminated without
raising, so no exception was captured and the proximate cause is not in the result
files. What the timing shows is that they are not startup failures and not a bad node.
No crash occurs before $1{,}063$ seconds; the median is $2{,}069$ and the maximum
$3{,}582$, just inside the limit. They are distributed evenly over the ten shards, at
between $32$ and $38$ each, so no single node or batch accounts for them. And $239$ of
the $352$ are Laporte5.

A run that survives at least eighteen minutes and is then killed without raising is
consistent with a resource limit reached during the search rather than a defect in the
model, but the campaign does not record enough to establish that: it would take the
scheduler's accounting record for the job, which is not part of the result files. The
obvious size explanation is separately excluded --- measured by $n\lvert T\rvert$ the
crashed runs are \emph{smaller} than the completed ones, median $400$ against $500$.

What is established either way is that the loss is not random with respect to
difficulty. On the instances LSS completed, SSPMF proves optimality $78\%$ of the time;
on the $352$ it lost, only $60\%$. The instances LSS never attempted are the harder
ones, so its totals are flattering and its rate is not comparable with the others.
Section~\ref{ssec:limits} records what follows from that.

\subsection{How far each method gets}\label{ssec:results-solved}

\emph{The question.} How many instances does each method close, and does the
decomposition compete with the compact models it was built to improve on?

This is the coarsest measurement in the section and the least explanatory. It
establishes a ranking, and that ranking is the starting point for everything after it,
but by itself it says nothing about \emph{why} the ranking is what it is.
Sections~\ref{ssec:results-bbc} and \ref{ssec:results-tight} take that up.

Table~\ref{tab:solves} reports, for each configuration and each family, the number of
instances solved to proven optimality against the number of runs that returned a
result.

\begin{table}[t]\centering
\scriptsize
\setlength{\tabcolsep}{4pt}
\begin{tabular}{@{}lrrrrrrr@{}}
\toprule
Configuration & Catanzaro & Crama & Laporte3 & Laporte4 & Laporte5 & Laporte7 & All\\
\midrule
SSPMF              & $73/171$ & $68/160$ & $340/340$ & $330/330$ & $161/336$ & $66/80$ & $1038/1417$\\
\texttt{CATZ-F4}   & $55/171$ & $44/160$ & $340/340$ & $330/330$ & $72/327$  & $44/78$ & $885/1406$\\
LSS                & $52/132$ & $42/124$ & $340/340$ & $330/330$ & $39/101$  & $41/42$ & $844/1069$\\
\midrule
\texttt{BBC-LP+T}  & $38/168$ & $34/156$ & $340/340$ & $192/330$ & $84/336$  & $39/80$ & $727/1410$\\
\texttt{BBC-LP}    & $38/169$ & $35/151$ & $340/340$ & $191/329$ & $83/338$  & $39/80$ & $726/1407$\\
\texttt{BBC-K}     & $39/166$ & $34/145$ & $340/340$ & $190/329$ & $83/335$  & $39/80$ & $725/1395$\\
\texttt{BBC-LP+F+H}& $43/171$ & $37/160$ & $316/340$ & $157/330$ & $97/339$  & $40/80$ & $690/1420$\\
\texttt{BBC-LP+F}  & $37/169$ & $32/160$ & $313/340$ & $156/329$ & $85/340$  & $36/80$ & $659/1418$\\
\texttt{BBC-LP+F+C}& $37/170$ & $32/160$ & $311/340$ & $157/329$ & $84/336$  & $36/79$ & $657/1414$\\
\texttt{BBC-LP+ACC}& $43/170$ & $37/160$ & $286/340$ & $154/330$ & $97/339$  & $39/80$ & $656/1419$\\
\texttt{BBC-K+F}   & $37/171$ & $21/121$ & $316/340$ & $156/329$ & $86/338$  & $35/79$ & $651/1378$\\
\texttt{BBC-LP+F+P}& $35/171$ & $32/160$ & $282/340$ & $155/328$ & $83/336$  & $35/79$ & $622/1414$\\
\bottomrule
\end{tabular}
\caption{Instances solved to proven optimality, against runs that returned a result,
under a one-hour limit. Denominators below the family size are the failed runs of
Section~\ref{ssec:validation}; the LSS row is the one where this matters, and its
totals should not be read against the others.}
\label{tab:solves}
\end{table}

Three things can be read from Table~\ref{tab:solves}.

The compact models dominate the decomposition. SSPMF closes $1{,}038$ instances against
$727$ for the best Benders configuration, and that gap is far larger than anything the
missing runs could account for.

Among the compact models themselves the ranking is less clear than the totals suggest,
because LSS did not complete a third of its runs. On the $1{,}067$ instances that LSS
and SSPMF both attempted, LSS proves $844$ optima and SSPMF $827$; on the $1{,}066$
that LSS and \texttt{CATZ-F4} both attempted, \texttt{CATZ-F4} proves $865$ against LSS's
$844$. Since the runs LSS lost are the harder ones, its position in both comparisons is
flattered by an unknown amount. The honest reading is that the three compact models are
close to one another on this collection and that separating them would need the lost
runs recovered; what the campaign does establish is that all three are ahead of the
methods built here.

The families separate the methods rather than rank them uniformly. Laporte3 is solved
completely by six of the twelve configurations. Laporte4 splits them sharply: the
three compact models solve all $330$, and the Benders solver solves between $154$ and
$192$. Catanzaro and Crama are hard for everything, and no method closes more than
half of either.

Among the Benders variants, the three that do not separate at fractional points
(\texttt{BBC-LP}, \texttt{BBC-LP+T}, \texttt{BBC-K}) are the three best, within two
instances of one another. Every configuration carrying \texttt{+F} is worse than all
three. Section~\ref{ssec:results-bbc} takes that apart.

\paragraph{Solving times.}
Table~\ref{tab:times} reports times on the $535$ instances that all eight of the listed
configurations solve. The shifted geometric mean with a shift of ten seconds is used,
following standard practice for exact-method comparisons, because it is insensitive to
both very small and very large individual values.

\begin{table}[h]\centering
\footnotesize
\begin{tabular}{@{}lrrr@{}}
\toprule
Method & shifted geometric mean (s) & median (s) & maximum (s)\\
\midrule
SSPMF              & $1.90$  & $0.140$  & $48.7$\\
\texttt{BBC-K}     & $14.50$ & $0.018$  & $845.8$\\
\texttt{BBC-LP+T}  & $16.03$ & $0.066$  & $1126.0$\\
\texttt{BBC-LP}    & $16.11$ & $0.067$  & $1128.4$\\
\texttt{BBC-LP+F+H}& $30.74$ & $0.003$  & $3584.6$\\
\texttt{BBC-LP+F}  & $33.73$ & $0.107$  & $3586.6$\\
LSS                & $37.68$ & $18.254$ & $2211.5$\\
\texttt{CATZ-F4}   & $92.83$ & $74.304$ & $3506.5$\\
\bottomrule
\end{tabular}
\caption{Solving times on the $535$ instances that all eight configurations solve. The
Benders solver has by far the lowest median and a high mean, which is the signature of
a method that either finishes almost at once or struggles.}
\label{tab:times}
\end{table}

The pattern in Table~\ref{tab:times} is as informative as the ranking. The Benders
solver's median is one to three orders of magnitude below every compact model's, while
its mean is an order of magnitude above SSPMF's. It either closes an instance
immediately or does not close it at all. Section~\ref{ssec:results-tight} identifies
what separates the two cases, and it is not the size of the instance that does so.

\paragraph{How the solved count accumulates.}
The shifted geometric mean summarises the instances a method solves; it says nothing
about the ones it does not. Table~\ref{tab:cumulative} gives the whole curve: how many
instances each method has closed by one second, ten seconds, and so on to the limit.

\begin{table}[h]\centering
\footnotesize
\begin{tabular}{@{}lrrrrrrr@{}}
\toprule
Method & $1$s & $10$s & $60$s & $300$s & $900$s & $3600$s & attempted\\
\midrule
SSPMF              & $482$ & $735$ & $946$ & $959$ & $982$ & $1038$ & $1417$\\
\texttt{CATZ-F4}   & $128$ & $254$ & $444$ & $662$ & $795$ & $885$  & $1406$\\
LSS                & $57$  & $246$ & $532$ & $759$ & $819$ & $844$  & $1069$\\
\midrule
\texttt{BBC-LP}    & $403$ & $469$ & $525$ & $572$ & $650$ & $726$  & $1407$\\
\texttt{BBC-K}     & $431$ & $490$ & $531$ & $570$ & $699$ & $725$  & $1395$\\
\texttt{BBC-LP+F}  & $362$ & $435$ & $466$ & $508$ & $542$ & $659$  & $1418$\\
\bottomrule
\end{tabular}
\caption{Instances closed within a given time. The two groups have different shapes.}
\label{tab:cumulative}
\end{table}

\begin{figure}[h]\centering
\begin{tikzpicture}
\begin{semilogxaxis}[reportaxis, width=0.95\linewidth, height=6.6cm,
  xlabel={wall-clock time (seconds, log scale)}, ylabel={instances closed},
  xmin=0.05, xmax=3600, ymin=0, ymax=1120,
  xtick={0.1,1,10,60,300,3600},
  xticklabels={$0.1$,$1$,$10$,$60$,$300$,$3600$},
  ytick={0,200,400,600,800,1000},
  yticklabel style={/pgf/number format/fixed,/pgf/number format/1000 sep={,}},
  legend pos=south east, legend columns=2, transpose legend]
  \addplot[black,    no marks]                  table[x=t,y=SSPMF]   {figdata/cactus.dat}; \addlegendentry{SSPMF}
  \addplot[black,    no marks, dashed]          table[x=t,y=CATZF4]  {figdata/cactus.dat}; \addlegendentry{\texttt{CATZ-F4}}
  \addplot[black,    no marks, dotted]          table[x=t,y=LSS]     {figdata/cactus.dat}; \addlegendentry{LSS}
  \addplot[black!55, no marks]                  table[x=t,y=BBCLP]   {figdata/cactus.dat}; \addlegendentry{\texttt{BBC-LP}}
  \addplot[black!55, no marks, dashed]          table[x=t,y=BBCK]    {figdata/cactus.dat}; \addlegendentry{\texttt{BBC-K}}
  \addplot[black!55, no marks, dotted]          table[x=t,y=BBCLPpF] {figdata/cactus.dat}; \addlegendentry{\texttt{BBC-LP+F}}
\end{semilogxaxis}
\end{tikzpicture}
\caption{Every solved run in the campaign, not the six sampled columns of
Table~\ref{tab:cumulative}: for each method, the number of instances closed by each
wall-clock time. The curves cross, and where they cross is what matters here. Before about
two seconds the configurations built here lead the field --- \texttt{BBC-K} has a
median solve time of $0.06$ seconds against $59.8$ for \texttt{CATZ-F4} --- because a
Benders master with a tight coverage row closes an instance at the root immediately or
not at all. After two seconds they stop moving and every compact model overtakes them.
The decomposition is not running out of time. It has stopped making progress, and what
would restart it is a bound that improves when the solver branches.}
\label{fig:cactus}
\end{figure}

The shapes differ, and the difference is the point. \texttt{BBC-LP} closes $403$
instances in the first second and $726$ in the full hour: more than half of everything
it will ever solve is solved before a compact model has finished building its model.
After that it flattens. SSPMF starts lower relative to its own total, at $482$ of
$1{,}038$, and keeps climbing throughout; LSS closes $57$ in the first second and
$844$ by the end. Extending the limit would therefore not affect the methods equally.
The Benders solver is not slow. It is fast on one class of instance and unable to
finish on the other, which is the behaviour Section~\ref{ssec:results-tight} explains.

\paragraph{How far from proven the unsolved runs are.}
The complement of the solved count is the gap remaining when the limit is reached,
measured between the incumbent and the dual bound. For \texttt{BBC-LP} the median over
its $681$ time-limited runs is $22.9\%$ and the mean is $28.6\%$, with a maximum of
$73\%$; the nine Benders configurations lie between $26.0\%$ and $29.5\%$ on the mean.
Only $50$ of those $681$ runs finish within five per cent of proven.

That figure should be read against the one in the next paragraph. On the runs where the
optimum is known independently, the solver's incumbent already equals it $59\%$ of the
time. It is therefore sitting on the right answer with a certified gap of roughly a
quarter, and the hour is spent failing to close a gap that exists only in the
certificate.

\subsection{The Benders solver and its strengthenings}\label{ssec:results-bbc}

\emph{The question.} A branch-and-bound method can fail for two reasons: it cannot find
a good solution, or it finds one and cannot prove it optimal. The two call for opposite
remedies, and the ablation below is designed to tell them apart. Four strengthenings
were implemented, one aimed at the incumbent and three at the dual bound. If the
difficulty is search, the first should be marginal and the others should help. If the
difficulty is the bound, the pattern should be reversed.

That is what makes this subsection informative rather than merely a list of settings:
the outcome discriminates between two explanations, and it was not obvious in advance
which one it would pick.

\paragraph{Fractional cuts.}
Separating Benders cuts at fractional points is a standard acceleration whose value on
this problem had not been established. The measurement below settles it.

The cuts are now generated in very large numbers. Across the $1{,}415$ instances that
both \texttt{BBC-LP} and \texttt{BBC-LP+F} attempted, the fractional configuration
added $67{,}513{,}396$ cuts, and the configuration without fractional separation added
none, which is the expected control.

They do not help. Table~\ref{tab:frac} gives the comparison.

\begin{table}[h]\centering
\footnotesize
\begin{tabular}{@{}lrr@{}}
\toprule
On the $1{,}415$ common instances & \texttt{BBC-LP} & \texttt{BBC-LP+F}\\
\midrule
fractional cuts added        & $0$   & $67{,}513{,}396$\\
solved to optimality         & $726$ & $659$\\
median nodes explored        & $77{,}504$ & $3{,}795$\\
\midrule
\multicolumn{3}{@{}l}{final dual bound, on the $1{,}404$ instances where both are recorded:}\\
\quad strictly higher with fractional cuts & \multicolumn{2}{r}{$0$}\\
\quad equal                                & \multicolumn{2}{r}{$1{,}228$}\\
\quad strictly lower with fractional cuts  & \multicolumn{2}{r}{$176$}\\
\bottomrule
\end{tabular}
\caption{The effect of enabling fractional Benders cuts. Sixty-seven million cuts
enter the master. The bound at termination is never higher than without them, and on
$176$ instances it is lower. The node counts are consistent with the hour going into
generating and managing cuts instead of searching, though the campaign does not isolate
that mechanism. Sixty-seven instances that the base solver closes are lost.}
\label{tab:frac}
\end{table}

The mechanism works:
the cuts are valid, they are generated, and they enter the master. The node counts show
they do tighten the relaxation locally, since the search explores twenty times fewer
nodes. What they do not do is raise the bound that matters. On no instance in the
comparison is the bound at termination higher with them than without.

The reading is that the fractional dual information is not new information. The
integer cuts already give the master everything the loading dual has to say, and the
fractional ones re-derive it at a cost of an hour of separation.

\paragraph{The three strengthenings.}
Table~\ref{tab:ablation} compares each strengthening against \texttt{BBC-LP+F} on the
instances both attempted.

\begin{table}[h]\centering
\footnotesize
\begin{tabular}{@{}llrrr@{}}
\toprule
Configuration & what it adds & common & solved & difference\\
\midrule
\texttt{BBC-LP+F}   & ---                        & ---       & ---   & ---\\
\texttt{BBC-LP+F+H} & hybrid-genetic incumbent   & $1{,}421$ & $690$ & $+31$\\
\texttt{BBC-LP+F+C} & conflict-graph row         & $1{,}421$ & $657$ & $-2$\\
\texttt{BBC-LP+ACC} & all three together         & $1{,}421$ & $656$ & $-3$\\
\texttt{BBC-LP+F+P} & Pareto-optimal cut lifting & $1{,}421$ & $622$ & $-37$\\
\midrule
\texttt{BBC-LP}     & fractional separation removed & $1{,}415$ & $726$ & $+67$\\
\bottomrule
\end{tabular}
\caption{Each strengthening against \texttt{BBC-LP+F}, on the instances both attempted.
Only the one that improves the incumbent helps. The two that aim at the dual bound
cost more time than they return, and combining all three does not recover the loss.
The last row is the comparison of Table~\ref{tab:frac}, repeated here because removing
a feature outperforms every addition.}
\label{tab:ablation}
\end{table}

The pattern is consistent and it is the same pattern as the fractional cuts. The only
strengthening that helps is the one that supplies a better \emph{incumbent}. Every
attempt to improve the \emph{bound} costs more than it returns. The conflict-graph row
is neutral, which is what Section~\ref{ssec:accel} predicted for it on structural
grounds: on most instances the coverage row already dominates it.

\paragraph{Where the search stops.}
Two further measurements locate the difficulty precisely.

First, of the $726$ instances \texttt{BBC-LP} solves, $259$ are closed at the root
node, without branching at all. The root relaxation value was recorded on $1{,}178$ of
the $1{,}421$ runs, and on $883$ of those the optimum is also known. On those $883$ the
root value equals the optimum $294$ times, and the split is sharp: among the runs that
went on to close the instance it is already the optimum in $264$ of $497$, and among
the runs that did not close it in only $30$ of $386$. Where it is not the optimum it is
short of it by at least one unit and by $16.7\%$ at the median, rising to $42.9\%$.

Second, consider the runs that hit the time limit. Across the nine Benders
configurations there are $6{,}562$ of them, and for $3{,}874$ the optimum is known
because some method proved it. On $2{,}571$ of those $3{,}874$, which is $66\%$, the
Benders solver's incumbent already \emph{equals} that optimum. It had found the answer
and could not prove it. For \texttt{BBC-LP} alone the figure is $229$ of $386$, or
$59\%$.

Taken together these two measurements describe a solver that finishes at the root when
the bound is good, and that finds the optimum quickly and then spends an hour failing to
certify it when the bound is bad. In neither case is the search the constraint on what
it can do.

\subsection{The coverage bound decides the outcome}\label{ssec:results-tight}

Section~\ref{sec:gap} predicted that the coverage bound $q$ would separate easy
instances from hard ones for any configuration-based method. That prediction is
directly testable, because for each instance with a known optimum it can be checked
whether the optimum equals the bound.

Of the $1{,}115$ instances whose optimum is known, $523$ have the coverage bound tight
and $592$ do not. The split is close to even, at $46.9\%$ tight, which means the
standard instance collections are not concentrated in either regime. Where the bound is
loose it is loose by a wide margin: the optimum exceeds it by $5.71$ on average and by
as much as $30$.

Splitting the results by that criterion gives Table~\ref{tab:tightness}.

\begin{table}[h]\centering
\footnotesize
\begin{tabular}{@{}lrrrr@{}}
\toprule
& instances & \texttt{BBC-LP} solves & SSPMF solves\\
\midrule
coverage bound tight ($\Zst=q$)  & $523$ & $492$ \; ($94\%$) & $523$ \; ($100\%$)\\
coverage bound loose ($\Zst>q$)  & $592$ & $234$ \; ($40\%$) & $515$ \; ($87\%$)\\
\bottomrule
\end{tabular}
\caption{Success rate split by whether the coverage bound equals the optimum. The
bound, rather than the size of the instance, is what the outcome tracks for the
Benders solver. The multicommodity model is affected by the same split but far less. A natural
reading is that a compact model gives its search a full description of the problem at
every node, so the search recovers what the bound does not supply; the campaign shows the
difference but does not test that explanation.}
\label{tab:tightness}
\end{table}

\begin{figure}[h]\centering
\begin{tikzpicture}
\begin{axis}[reportaxis, height=6.2cm,
  xlabel={how loose the coverage bound is on the instance \ (\,$\Zst-q$\,)},
  ylabel={instances closed (\%)},
  xmin=-0.6, xmax=8.6, ymin=0, ymax=112,
  xtick={0,1,2,3,4,5,6,7,8},
  xticklabels={$0$,$1$,$2$,$3$,$4$,$5$--$6$,$7$--$9$,$10$--$14$,$15+$},
  ytick={0,20,40,60,80,100},
  legend style={at={(0.03,0.06)}, anchor=south west}]
  \addplot[black, mark=*, mark size=1.9]
    coordinates {(0,100.0)(1,100.0)(2,100.0)(3,100.0)(4,100.0)(5,96.6)(6,84.3)(7,63.9)(8,9.3)};
  \addlegendentry{SSPMF (compact)}
  \addplot[black!55, mark=o, mark size=1.9, dashed]
    coordinates {(0,94.1)(1,30.4)(2,37.8)(3,42.9)(4,60.7)(5,34.1)(6,47.0)(7,52.5)(8,9.3)};
  \addlegendentry{\texttt{BBC-LP} (this report)}
  \node[font=\scriptsize, text=black!55, anchor=west] at (axis cs:0.15,104) {$523$ instances};
  \node[font=\scriptsize, text=black!55, anchor=east] at (axis cs:7.85,15) {$43$ instances};
\end{axis}
\end{tikzpicture}
\caption{The $1{,}115$ instances of known optimum, grouped by how far the optimum
exceeds the coverage bound, against the share each method closes. This is
Table~\ref{tab:tightness} resolved into its whole response curve, and the curve says
more than the two-way split does. The Benders solver does not degrade: it falls off a
cliff between a bound that is exact and a bound that is wrong by one, from $94\%$ to
$30\%$, and then wanders without trend --- once the bound is loose at all, by how much
stops mattering. The compact model degrades smoothly instead, holding $100\%$ out to a
gap of four. But the two meet again at the right-hand end: on the $43$ instances whose
bound is loose by fifteen or more, \emph{both} close $9.3\%$. The compact model is not
immune to the bound; it has further to fall.}
\label{fig:tightsplit}
\end{figure}

This is the central measurement of the campaign. The same solver, on the same hardware,
under the same time limit, closes $94\%$ of one class and $40\%$ of the other. The
classification is not by number of jobs, number of tools or density; it is by a single
structural property of the instance, and it is the property the theory of
Section~\ref{sec:gap} identifies.

The SSPMF column is the useful control. Its relaxation is also equal to $q$, so if the
bound were the whole story it should show the same collapse. It does not: it falls from
$100\%$ to $87\%$ where the Benders solver falls from $94\%$ to $40\%$. A compact model
gives the branch-and-bound search a full description of the problem to work with at
every node, and the search recovers much of what the bound does not supply. The Benders
master sees only the cuts it has accumulated, and when the bound is loose it has very
little else.

\subsection{Branch-and-price: the relaxation measured directly}\label{ssec:results-bnp}

\emph{The question.} Everything above measures the bound indirectly, through how many
instances a solver closes. The prototypes can answer directly, because they report the
value of their root relaxation on every run. Two propositions of Section~\ref{ssec:pos}
predict what that value will be --- that PCF$'$ equals the coverage bound and that PTF
can exceed it --- and the campaign is a chance to test a proof against a benchmark
rather than against another proof.

That is the whole purpose of this run, and it is worth being explicit about what it is
not. It is not a fourth entry in the comparison of Section~\ref{ssec:results-solved}.
The prototypes are implemented in Python over PySCIPOpt with a Python pricer, so a
single node costs orders of magnitude more than a node of a compiled solver; their
wall-clock times measure the implementation language, not the formulation, and are
never compared against CPLEX anywhere in this report. The quantity this subsection
exists to report is the root relaxation value, which is a property of the formulation
and is recorded before any search happens. The solve counts are given because a
reader will ask for them, and are read only as evidence that the prototypes are
working code.

In total $2{,}581$ runs were made over $469$ instances in seven
configurations: PCF$'$ alone; PCF$'$ with each of the four accelerations of
Section~\ref{ssec:pos} separately, written \texttt{+HP} for heuristic pricing,
\texttt{+MC} for multiple pricing, \texttt{+WS} for warm starting and \texttt{+STAB}
for dual stabilisation; PCF$'$ with all four together, written \texttt{+ACC}; and PTF.
The time limit is thirty minutes on the Catanzaro and Crama families and ten minutes on
the Laporte families.

Two properties of this run make its counts weaker still, and both are stated before the
numbers. First, the harness carries an early-stop rule: once a configuration has failed
to close six increasingly hard instances in succession, the remaining harder instances
are skipped for that configuration. The number of runs a configuration received
therefore depends on how it performed, so the denominators in Table~\ref{tab:bnp} are
not comparable across rows. Second, no repetition or seed variation was run. Every
comparison drawn below is made either on instances that both configurations attempted,
or on the root relaxation value, which the early-stop rule cannot bias because it is
recorded before any search.

\paragraph{They agree with the campaign.}
On the $507$ runs that proved an optimum and where the same instance was also closed in
the campaign of Section~\ref{ssec:results-solved}, the two agree on every one. As in
Section~\ref{ssec:validation}, the native objective is the free-initial cost and the
difference from the empty-start cost is exactly $b$ on all $507$.

\paragraph{What the root relaxation shows.}
Proposition~\ref{prop:pcf} states that the position--configuration relaxation equals
the coverage bound $q$, and Proposition~\ref{prop:ptf} states that the
position--transition relaxation can exceed it. Both are testable directly, because the
prototypes record the root relaxation value.

\emph{PCF$'$ never exceeds $q$.} On all $2{,}114$ PCF$'$ runs for which a root
relaxation value was recorded, that value equals $\lvert\U\rvert-b$ exactly. Not one
run, in any of the six PCF$'$ configurations, produced a root bound above the counting
bound. This is Proposition~\ref{prop:pcf} confirmed on the benchmark rather than on
paper.

\emph{PTF exceeds $q$, and almost never.} Of the $233$ PTF runs with a recorded root
relaxation, three are strictly above $\lvert\U\rvert-b$: Catanzaro \texttt{A0-0}, where
the bound is $2.013$ against $q=2$, and Laporte3 \texttt{L11-6} and \texttt{L11-7},
where it is $6$ against $q=5$. Proposition~\ref{prop:ptf} is therefore not vacuous. But
$230$ of $233$ runs sit exactly on $q$, and on the two Laporte3 instances the optima
are $12$ and $13$, so a bound of $6$ instead of $5$ closes a small part of a large gap.

\emph{The bound is often already the answer.} On $990$ of the $1{,}846$ runs where both
the root relaxation and the optimum are known, the root relaxation \emph{equals} the
optimum. These are the tight instances of Section~\ref{ssec:results-tight}, seen from
the other side: where the coverage bound is tight, a configuration-based relaxation
needs no strengthening at all, and where it is loose, neither PCF$'$ nor PTF supplies
one.

\paragraph{What the prototypes close, and why that is not the measurement.}
Table~\ref{tab:bnp} gives the counts, for completeness rather than for comparison.
Nothing above nine jobs was closed by any configuration: all $608$ runs on the $30$-job
and $40$-job instances reached the time limit. This is the expected consequence of a
Python pricer and is not evidence about the formulation, which the paragraph above has
already measured.

\begin{table}[h]\centering
\footnotesize
\begin{tabular}{@{}llrrrr@{}}
\toprule
Configuration & what it adds & runs & solved & median time (s) & median columns\\
\midrule
PCF$'$\texttt{+MC}   & multiple pricing          & $422$ & $102$ & $11.4$ & $1{,}504$\\
PCF$'$               & ---                       & $415$ & $91$  & $23.7$ & $1{,}574$\\
PCF$'$\texttt{+HP}   & heuristic pricing         & $360$ & $73$  & $19.7$ & $5{,}173$\\
PTF                  & position--transition rows & $387$ & $73$  & $2.4$  & $1{,}279$\\
PCF$'$\texttt{+WS}   & warm-started pool         & $359$ & $64$  & $14.9$ & $1{,}610$\\
PCF$'$\texttt{+ACC}  & all four together         & $337$ & $60$  & $28.1$ & $1{,}221$\\
PCF$'$\texttt{+STAB} & dual stabilisation        & $301$ & $44$  & $35.7$ & $1{,}133$\\
\bottomrule
\end{tabular}
\caption{The branch-and-price prototypes, reported as evidence that the code runs and
not as a performance comparison. Median time is over the runs that closed; median
columns is over all runs. The run counts are outcome-dependent through the early-stop
rule, so the rows are not comparable with one another, and none of them is comparable
with the compiled solvers of Table~\ref{tab:solves}. The measurement this subsection
exists for is the root relaxation above.}
\label{tab:bnp}
\end{table}

Head to head on the $320$ instances both attempted, PCF$'$ and PTF each close $73$.
Their node counts differ by a factor of fifty --- a median of $96$ against $2$ ---
because the PTF pricer prices coupled pairs of configurations and costs far more per
node. The stronger relaxation does not convert into more instances closed.

\paragraph{The comparison with the Job Grouping Problem.}
\citet{Otiai2024} builds a branch-and-price for that problem whose set-covering
relaxation is strong enough to close benchmark instances in a handful of nodes. For the
SSP, \citet{Moreira2026} reports that the trivial bound $\lvert T\rvert-b$ obtained by
the multicommodity model was the best available bound on all thirty instances of three
of his subgroups, with the configuration model giving the best bound on four instances
of the remaining, tighter subgroup. The measurements above say the same thing on a
larger collection and with the mechanism visible: the configuration relaxation for the
SSP is not, in general, an improvement on the counting bound.

\subsection{What the campaign establishes}\label{ssec:results-summary}

Two of the results above are checks that the study is sound. Cross-method agreement is
complete: on the $999$ instances solved by more than one method, all $41{,}673$ pairs
agree on the cost of the schedule returned. And the three reimplemented compact models
all close substantially more instances than any configuration built here, without being
cleanly separable from one another, because the LSS runs that failed are the harder
ones.

The remaining four converge on one statement. Whether the coverage bound equals the
optimum decides whether the Benders solver closes an instance, and where it fails the
incumbent is usually already optimal. Of four strengthenings, only the one that
improves the incumbent helps; the three aimed at the bound are neutral or harmful, and
the sixty-seven million fractional cuts are the extreme case, raising the bound at
termination on no instance while costing sixty-seven solved ones. Both branch-and-price
relaxations sit on the coverage bound or barely above it, confirming
Propositions~\ref{prop:pcf} and \ref{prop:ptf} and showing the second to be of little
practical use as it stands. Section~\ref{sec:disc} takes that up.
